% Standalone companion; original Revision 11 and production code are preserved.
\documentclass[11pt]{article}
\usepackage[T1]{fontenc}
\usepackage[utf8]{inputenc}
\usepackage{amsmath,amssymb,mathtools,booktabs,array,longtable}
\usepackage[margin=1in]{geometry}
\usepackage[hidelinks]{hyperref}
\setlength{\emergencystretch}{3em}
\newcommand{\kb}{k_{\rm B}}
\newcommand{\dd}{\mathrm d}
\newcommand{\E}{\mathbb E}
\newcommand{\Pm}{\mathcal P}
\newcommand{\Qm}{\mathcal Q}
\title{From the Poisson--Bessel Lattice Energy to Reduced Dynamics:\
PART I--III Foundations and the Two Missing Bridges}
\author{Research companion to the Revision 11 working manuscript}
\date{20 October 2026}
\begin{document}
\maketitle
\begin{abstract}
We reconstruct the ideal lattice energy without discarding its nonzero
Poisson--Bessel modes, then distinguish it from the finite-temperature
potential of mean force and from a probability-current law. Explicit
collective coordinates, mass and measure give the conditional partition
integral. Conditional projection gives a generally nonlinear memory
equation. An exactly eliminable harmonic bath keeps the full nonlinear
Bessel potential while exhibiting Gaussian but recurrent colored forcing.
Separate Markov and small-mass assumptions yield a conditional diffusion,
including its noise-induced drift for variable mobility. Eight independent
mathematical checks support the identities and counterexamples in declared
dimensionless models. They do not calibrate aluminium or establish a
fatigue mechanism. Revision 11 is retained as a conditional framework;
same-coordinate thermodynamic, kinetic and observation gates remain open.
\end{abstract}
\begin{quote}\small
\textbf{Scope.} First delivery: PART I--III only. The full later mechanism,
spatial coupling, 55-display downstream audit and experimental fatigue-life
validation are deferred. Original Revision 11, production solver,
calibration and historical results are preserved. This is a foundations
audit, not a substitution of hypothetical numbers for missing material data.
\end{quote}
\tableofcontents

\section{The two questions}
The mathematical identity for an ideal rigid lattice is not the disputed
step. The two questions are
\begin{equation}
W_\Gamma^{\rm Bessel}\ \xrightarrow[\text{same coordinate and measure}]
 {\text{integrate out atoms}}\ A(q)
\ \xrightarrow[\text{time scales and preparation}]
 {\text{reduce actual dynamics}}\ \boldsymbol J.
\end{equation}
Here $q=(\delta a,s)$, with normal cut opening and tangential registry.
An additive shift to $a$ changes no derivative. $A$ is a free energy,
not a cut area. Geometry, coherent patch size, temperature, load, initial
ensemble and boundaries must be the same through both arrows.

\section{PART I: lattice energy, counting and derivatives}
\subsection{Hamiltonian and the row transform}
For finite atom positions $R$ and momenta $\Pi$,
\begin{equation}
H=\tfrac12\Pi^Tm^{-1}\Pi+V(R),\qquad
V_{\rm pair}=\sum_{i<j}4\epsilon
 [(\sigma/r_{ij})^{12}-(\sigma/r_{ij})^6].
\end{equation}
The pair force is $-\phi'(r)\widehat r$ and its radial Hessian is
$\phi''(r)\widehat r\widehat r^T+
(\phi'(r)/r)(I-\widehat r\widehat r^T)$, with
\begin{align}
\phi'&=4\epsilon[-12\sigma^{12}r^{-13}+6\sigma^6r^{-7}],\\
\phi''&=4\epsilon[156\sigma^{12}r^{-14}-42\sigma^6r^{-8}].
\end{align}
For $a>0$, $p>1/2$, define $x_j=(j+1/2)b-s$ and
\begin{equation}
S_p(a,s)=\sum_{j\in\mathbb Z}(a^2+x_j^2)^{-p},\qquad
W_{\rm row}=4\epsilon(\sigma^{12}S_6-\sigma^6S_3).
\end{equation}
Using $z^{-p}=\Gamma(p)^{-1}\int_0^\infty t^{p-1}e^{-tz}\dd t$,
the Gaussian Fourier integral gives, for $k\ne0$,
\begin{align}
\widehat f_p(k)&=\frac{\sqrt\pi}{\Gamma(p)}
 \int_0^\infty t^{p-3/2}e^{-a^2t-k^2/(4t)}\dd t\\
 &=\frac{2\sqrt\pi}{\Gamma(p)}
 \left(\frac{|k|}{2a}\right)^{p-1/2}K_{p-1/2}(a|k|),\\
\widehat f_p(0)&=\frac{\sqrt\pi\Gamma(p-1/2)}{\Gamma(p)}a^{1-2p}.
\end{align}
The last integral is the modified-Bessel representation\cite{DLMF}.
Poisson summation then gives
\begin{align}
S_p={}&\frac{\sqrt\pi\Gamma(p-1/2)}{b\Gamma(p)}a^{1-2p}\nonumber\\
&+\frac{4\sqrt\pi}{b\Gamma(p)}\sum_{n=1}^\infty
 \left(\frac{\pi n}{ab}\right)^{p-1/2}
 K_{p-1/2}(2\pi na/b)\cos[2\pi n(1/2-s/b)].
\label{eq:row}
\end{align}
Thus LJ rows use $K_{5/2}$ and $K_{11/2}$. The zero reciprocal mode is a
lateral average, not a bulk environment. All nonzero modes are retained.
On compact sets bounded away from $a=0$, the sums and their first two
derivative sums converge uniformly. Direct differentiation gives
\begin{align}
S_{p,a}&=-2paS_{p+1},& S_{p,s}&=2p\sum_j x_j(a^2+x_j^2)^{-p-1},\\
S_{p,aa}&=-2pS_{p+1}+4p(p+1)a^2S_{p+2},\nonumber\\
S_{p,as}&=-4p(p+1)a\sum_jx_j(a^2+x_j^2)^{-p-2},\nonumber\\
S_{p,ss}&=-2pS_{p+1}+4p(p+1)\sum_jx_j^2(a^2+x_j^2)^{-p-2}.
\end{align}

\subsection{An explicit cut geometry and its multiplicities}
Choose the illustrative zigzag stack
$r_{j\ell}=(jb+\ell b/2,\ell h)$. Lower layers have $\ell\leq0$;
upper layers $\ell\geq1$ translate by $(-s,\delta a)$. With layer difference
$k=\ell_+-\ell_-$, the correct row-kernel argument is
\begin{equation}
w_k(\delta a,s)=W_{\rm row}(kh+\delta a,s+(1-k)b/2).
\end{equation}
There are exactly $k$ crossing layer pairs for each $k\geq1$, since
$1-k\leq\ell_-\leq0$. Hence, per lateral primitive cell,
\begin{equation}
W_\Gamma(\delta a,s)=\sum_{k\geq1}k
 [w_k(\delta a,s)-w_k(0,0)].\label{eq:cut}
\end{equation}
For $L$ lower and $U$ upper layers the count is instead
\begin{equation}
c_k=\max\{0,\min(0,U-k)-\max(1-L,1-k)+1\}.
\end{equation}
No further pair factor $2$ or $1/2$ applies. Crucially,
$\partial_{\delta a}(kh+\delta a)=1$, so
\begin{equation}
\partial_{\delta a}W_\Gamma=\sum_k k\,\partial_dw_k,
\qquad \partial_{\delta a}^2W_\Gamma=\sum_k k\,\partial_d^2w_k.
\end{equation}
The extra chain factor $k$ for a homogeneous dilation $d_k=ka$ is absent.
The absolute row baseline converges for $p>3/2$, including LJ. On
$\delta a>-h$, local uniform derivative convergence holds. For opening
tending to infinity at fixed registry, dominated convergence gives
\begin{equation}
W_{\rm sep}=-\sum_{k\geq1}k\,w_k(0,0).
\end{equation}
The zero mode, all registries and the complete crossing sum enter this
limit. Matched equivalent surfaces give $W_{\rm sep}=2A_\Gamma\gamma$;
rigid cleavage and relaxed surface formation are different targets.

\subsection{A full plane is a different transform}
For a planar Bravais lattice with primitive area $A_\parallel$,
\begin{align}
S_p^{(2)}(d,\Delta)&=A_\parallel^{-1}
 \sum_{G\in\Lambda^*}\widehat f_p^{(2)}(G)e^{iG\cdot\Delta},\\
\widehat f_p^{(2)}(G)&=\frac{2\pi}{\Gamma(p)}
 (|G|/(2d))^{p-1}K_{p-1}(d|G|),\quad G\ne0,\\
\widehat f_p^{(2)}(0)&=\frac{\pi}{p-1}d^{2-2p}.
\end{align}
Here $d>0,p>1$. LJ planes use $K_2,K_5$. The counted plane baseline
converges for $p>2$. For a declared FCC(111) example one can use triangular
vectors $a_1=b(1,0)$, $a_2=b(1/2,\sqrt3/2)$,
$b=a_{\rm lat}/\sqrt2$, $h=a_{\rm lat}/\sqrt3$ and
$\tau=(a_1+a_2)/3$, evaluating the plane kernel at
$(kh+\delta a,k\tau+u)$ with count $k$. Orientation, scalar slip path
and translation period remain required inputs. A row certificate is not
automatically a full-plane or counted-cut certificate.

For an embedding potential, first form each complete site density and then
sum $\sum_i[F(\rho_i(q))-F(\rho_i^{\rm ref})]$. Multiplying an individual
layer's embedding function by a pair-crossing count is not valid.

\section{PART II: atoms, PMF and exact memory}
\subsection{Collective coordinate, mass and hidden variables}
Take a finite, normalizable canonical reference, removing translation zero
modes consistently. Define $q=CR+q_c$ with rank-two constant $C$ and
\begin{equation}
G=Cm^{-1}C^T,\qquad M_q=G^{-1},\qquad B=m^{-1}C^TM_q.
\end{equation}
Choose $CE=0$ so $R=R_0+Bq+Ey$ and $B^TmE=0$. Then the canonical momentum
$p=B^T\Pi=M_q\dot q$ separates a kinetic energy $p^TM_q^{-1}p/2$.
For normal/tangent projections of two finite group mass-centre differences,
$M_q=\mu I_2$, $\mu=M_+M_-/(M_++M_-)$. Other observables have other masses.
The hidden variables include internal vibrations, common motion, unretained
relative motion and all allowed rearrangements. They are not automatically
fast. Infinite half-crystal masses diverge; a finite coherent patch has
energy $N_cW_\Gamma$ and mass proportional to $N_c$. A per-cell potential
cannot be combined with an unspecified whole-half-crystal mass.

\subsection{The free-energy bridge}
Relative to flat $\dd q$, define
\begin{equation}
Z(q)=\int\delta(\Xi(R)-q)e^{-\beta V(R)}\dd R,\qquad
A=-\kb T\log Z+C_A.\label{eq:pmf}
\end{equation}
The affine map has constant Jacobian. With
$V(q,y)=U_{\rm Bessel}(q)+V_\perp(q,y)$,
$Z_\perp=\int e^{-\beta V_\perp}\dd y$, one obtains exactly
\begin{equation}
\boxed{\Delta A=\Delta U_{\rm Bessel}
                  -\kb T\log[Z_\perp(q)/Z_\perp(q_0)].}\label{eq:bridge}
\end{equation}
Thus $A=U_{\rm Bessel}+C_A$ if and only if $Z_\perp$ is constant on the
adopted domain and measure. The Bessel identity proves no such condition.
The reference units inside logarithms are fixed; differences use ratios.

For nonlinear $\Xi$, coarea gives
\begin{equation}
Z(q)=\int_{\Xi^{-1}(q)}\frac{e^{-\beta V}}{J_\Xi}\dd\sigma_q,
\quad J_\Xi=\sqrt{\det(D\Xi D\Xi^T)}.
\end{equation}
If $b_i\cdot\nabla\Xi_j=\delta_{ij}$, integration by parts gives
\begin{equation}
\partial_iA=\langle b_i\cdot\nabla V-\kb T\nabla\cdot b_i\rangle_q,
\end{equation}
provided boundary terms vanish. The coordinate measure and constraint
ensemble determine metric corrections; do not append a generic correction
without those choices\cite{Legoll2010}. A variable mass metric along one
level set can give a mixture of conditional velocity Gaussians, rather
than a single exact inertial mass.

At a unique interior nondegenerate constrained minimum $y_*(q)$, use the
measure $w(q,y)\dd y$ and positive Hessian $H_\perp$. Laplace expansion gives
\begin{align}
A(q)={}&U_{\rm rel}(q)+\tfrac12\kb T\log\det H_\perp(q)
 -\kb T\log w(q,y_*)\nonumber\\
 &-\tfrac d2\kb T\log(2\pi\kb T)+C_A+O(T^2).
\end{align}
Ratios fix units in differences. Zero modes, multiple minima, boundary
minima and nonuniform limits need separate treatment. For the exact
counterexample $V=U_{\rm Bessel}(q)+\kappa(q)y^2/2$, $\kappa>0$,
\begin{equation}
U_{\rm rel}=U_{\rm rigid},\qquad
\Delta A=\Delta U_{\rm Bessel}+\tfrac12\kb T\log[\kappa(q)/\kappa(q_0)].
\end{equation}
The nonlinear Bessel potential has not been replaced by a parabola; a hidden
harmonic degree of freedom changes its PMF even though it leaves the
constrained minimum energy unchanged.

\subsection{Liouville marginal and conditional projection}
For any initial ensemble, the exact marginal continuity law is
\begin{equation}
\partial_tP=-\nabla_q\cdot J,\qquad J=P\E_{\rho_t}[\dot q\mid q].
\end{equation}
The evolving conditional hidden distribution prevents closure by this
identity alone. Smooth finite Hamiltonian positions have zero quadratic
variation; Brownian positions require a limiting description.

For an autonomous invariant canonical reference let $Z=(q,p)$,
$Lf=\{f,H\}$, $\Pm f=\E_{\rm eq}[f\mid Z]$, $\Qm=I-\Pm$.
On a suitable common operator domain the Dyson identity gives
\begin{equation}
\dot Z(t)=e^{tL}\Pm LZ+e^{t\Qm L}\Qm LZ
 +\int_0^t e^{(t-u)L}\Pm Le^{u\Qm L}\Qm LZ\,\dd u.
\label{eq:mz}
\end{equation}
For the affine coordinate, $\Qm Lq=0$, $\Pm Lp=-\nabla A$.
Put $F(u)=e^{u\Qm L}\Qm Lp$ and
$C_{ij}(z,u)=\E[F_i(u)F_j(0)\mid Z=z]$. Orthogonality gives $\Pm F=0$,
not white or independent force. Anti-self-adjointness of $L$ and integration
by parts in $p$ give
\begin{equation}
\boxed{\Pm LF_i(u)=\rho_Z^{-1}\partial_{p_j}(\rho_Z C_{ij})
 =\partial_{p_j}C_{ij}-\beta C_{ij}(M_q^{-1}p)_j.}
\end{equation}
For example, pairing with $g(Z)$ gives
$\langle g,LF_i\rangle=-\int\rho_Z\partial_{p_j}g\,C_{ij}\dd z$;
conditional orthogonality cancels the resolved part of $Lg$.
The memory in \eqref{eq:mz} evaluates this state-dependent expression at
$Z(t-u)$. A velocity convolution follows only in special cases, e.g.
$C$ independent of $z$. Conditional Zwanzig and linear Mori projections
must not be mixed without justification\cite{Mori1965,Zwanzig1973,Ayaz2022}.
No operator-domain theorem for every infinite LJ lattice is claimed here.

\subsection{Exactly eliminable bath atop the full Bessel energy}
Consider the specified finite Hamiltonian
\begin{equation}
H=\tfrac12p^TM_q^{-1}p+U(q)+\sum_j\left[
 \frac{P_j^2}{2m_j}+\frac{m_j\omega_j^2}{2}
 \left(x_j-\frac{c_j^Tq}{m_j\omega_j^2}\right)^2\right].
\end{equation}
Here $U$ is the full nonlinear Bessel potential. The shifted Gaussian bath
integral is independent of $q$, hence $A=U+C_A$ exactly. Solving each
oscillator and integrating by parts gives
\begin{align}
M_q\ddot q&=-\nabla U-\int_0^t\Gamma(t-u)\dot q(u)\dd u
                   +\eta(t)+F_{\rm ext}(t),\\
\Gamma(t)&=\sum_j\frac{c_jc_j^T}{m_j\omega_j^2}\cos\omega_jt,\\
\eta(t)&=\sum_jc_j\left[
 (x_j(0)-c_j^Tq(0)/(m_j\omega_j^2))\cos\omega_jt
 +\frac{P_j(0)}{m_j\omega_j}\sin\omega_jt\right].
\end{align}
Shifted conditional canonical preparation gives
\begin{equation}
\E\eta=0,\qquad \E[\eta(t)\eta(u)^T]=\kb T\Gamma(t-u).
\end{equation}
Finitely many modes yield Gaussian but colored, recurrent forcing.
Unshifted preparation instead gives an initial slip $-\Gamma(t)q(0)$.
Changing bath parameters changes kinetics while leaving the same $A$.
Off-diagonal memory requires symmetry or evidence to vanish. Arbitrary
cyclic forcing of $q$ leaves this special linear bath identity intact;
a driven anharmonic solid can instead require a nonstationary, two-time
memory and a preparation-dependent fluctuation relation.

\section{PART III: conditional stochastic limits and current}
\subsection{White-bath limit and Kramers dynamics}
A chosen exponentially decorrelating family has
\begin{equation}
\Gamma_\varepsilon(t)=\zeta\varepsilon^{-1}e^{-t/\varepsilon},\ t\geq0,
\quad\int_0^\infty\Gamma_\varepsilon\dd t=\zeta.
\end{equation}
Its stationary noise covariance
$\kb T\zeta\varepsilon^{-1}e^{-|t-u|/\varepsilon}$ tends to
$2\kb T\zeta\delta(t-u)$. This is a controlled bath family, not a spectral
derivation for the actual lattice. General non-Gaussian atomic baths need
mixing and a functional limit plus feedback/preparation control.
With constant symmetric $\zeta>0$, constant $T$ and constant mass, the
candidate Markov inertial equation is
\begin{align}
\dd q&=M_q^{-1}p\dd t,\\
\dd p&=[-\nabla A+F_{\rm ext}-\zeta M_q^{-1}p]\dd t
                         +\sqrt{2\kb T\zeta}\dd B_t.
\end{align}
Its Kramers equation retains momentum transport and momentum diffusion.
Integrating it over $p$ gives a current moment, not yet a position-only
Fick law. The conditions are separately
\begin{equation}
\tau_{\rm bath}\ll\tau_q,\tau_{\rm load},\qquad
\tau_p\ll\tau_q,\tau_{\rm load},\quad\tau_p=\mu/\zeta
\quad\hbox{(scalar case)}.
\end{equation}
The first supports Markov approximation; the second supports overdamping.
Joint colored-noise and small-mass limits require their ratio/order to be
declared. Absorbing boundaries can create additional inertial layers.

\subsection{Variable mobility and the forward equation}
For constant temperature and a smooth symmetric positive mobility $M(q)$,
the controlled small-mass limit of an inertial SDE with
$\sigma\sigma^T=2\kb T\zeta$ and $M=\zeta^{-1}$ is, in It\^o form,
\begin{equation}
\boxed{\dd q=[-M\nabla A+MF_{\rm ext}+\kb T\nabla\cdot M]\dd t
                  +B\dd B_t,\qquad BB^T=2\kb T M.}
\label{eq:ito}
\end{equation}
Regularity, positive friction, growth/nonexplosion or localization, initial
data and boundary assumptions remain required\cite{Hottovy2015}. This
theorem starts from an inertial SDE, not from unspecified atomic noise.
The It\^o forward equation gives
\begin{equation}
J_i=b_iP-\kb T\partial_j(M_{ij}P),\qquad
\boxed{J=-M[(\nabla A-F_{\rm ext})P+\kb T\nabla P].}
\label{eq:current}
\end{equation}
Omitting the mobility divergence in \eqref{eq:ito} changes this current.
A Stratonovich representation subtracts
$\tfrac12\sum_{jk}B_{jk}\partial_j B_{ik}$ from the It\^o drift; in
multiple dimensions that correction depends on the noise frame and is
not universally half the mobility divergence. Constant $M,T$ recover
Rev.11's algebra only after all earlier identification assumptions hold.

At fixed conservative external potential and admissible no-flux/decay
boundaries, a normalizable Gibbs law is stationary. Integration by parts
gives a nonpositive free-energy derivative
$-\int P[\nabla(A+V_{\rm ext})+\kb T\nabla\log P]^T
M[\nabla(A+V_{\rm ext})+\kb T\nabla\log P]\dd q$.
Absorbing populations and time-dependent work require their own extra terms.

\subsection{Domain and physical normalization}
The ideal rigid cut tends to finite energy at infinite opening and is
periodic in unwrapped slip. Its global Gibbs integral on that domain
diverges even without tensile load. Equilibrium projection needs a declared
finite confining or restricted reference; a periodic quotient loses winding
history unless an index is retained. A numerical reflecting box does not
prove physical confinement or a global continuous-law contract.

For pure normal traction,
$V_{\rm ext}=-A_\Gamma T_n(t)\delta a$, so $F_s=0$.
The old projected work $-f[(a-a_0)+\chi s]$ has a tangential component
$\chi f$ and requires its own loading geometry. Static energy does not
identify a mobility. With $q=\ell q^*$, $A=E_0A^*$, $M=M_0M^*$,
\begin{equation}
t_0=\ell^2/(M_0E_0),\qquad\theta=\kb T/E_0.
\end{equation}
A fully specified physical Hamiltonian has Newtonian time; the missing
quantity from an energy-only overdamped model is its compatible bath and
mobility. Dividing a coherent plane energy by $N_p$ requires both
$M\mapsto N_pM$ and $\kb T\mapsto\kb T/N_p$ to preserve the generator.
Multiplying mobility alone is not equivalent.

Existing repository v31 studies include actual Al99 MD and finite-band
plane-coordinate kinetic estimates. Their error sets include zero drag,
and their plane PMF/local-coordinate and thermal normalization match is
unresolved. They are evidence, but do not presently calibrate this local
physical clock. No measured seconds, hertz or fatigue cycles are inferred.

\section{Executed mathematical checks}
All parameters below are illustrative and dimensionless. The independent
script saved all truncations, tolerances and channel errors; it ran to
exit code zero with eight passes and no failures. No production regression
or application suite is claimed.
\begin{center}\small
\begin{tabular}{p{0.34\linewidth}p{0.58\linewidth}}
\toprule
Check & Evidence and scope\\\midrule
Row direct versus Poisson & Maximum scaled value/gradient/Hessian error
$9.84\,10^{-15}$; near-singular large absolute values are not assigned that
absolute error.\\
Triangular plane transform & Real/reciprocal cutoffs refined; recorded
analytic tail bound covers energy, not a Hessian tail certificate.\\
Finite crossing counts & $4\times5$ layers, 20 crossing row pairs;
direct/counted difference $9.44\,10^{-16}$.\\
Cut derivatives and endpoint & Stress-free stack $h=0.761436142462363$;
positive Hessian eigenvalues $60.8604,141.4438$;
$W_{\rm sep}=2.563620429191983$.\\
Exact PMF counterexample & At $(0.1,0.17)$, $\kb T=0.2$,
$\Delta U=0.5170967738516374$, $\Delta A=0.5458266268875693$;
same Bessel potential, hidden stiffness entropy.\\
Finite Hamiltonian bath & Refined nonlinear force identity residual
$1.23\,10^{-15}$; external-work balance $1.04\,10^{-14}$;
analytic covariance, recurrence and initial slip checked.\\
Two limit controls & Explicit local Bessel-Hessian OU model: equilibrium
covariance residual $8.00\,10^{-15}$; separate memory and mass refinements
reduce frequency-response errors.\\
Variable-mobility drift & On a fixed-opening periodic quotient, Gibbs flux
residual $3.56\,10^{-15}$ with correct drift; $2.4654$ without it.\\
\bottomrule
\end{tabular}
\end{center}
Two truncations agreeing at machine precision do not prove zero omitted
tail. The local quadratic limit control is explicitly an approximation;
the finite-bath force check keeps the nonlinear Bessel energy. No Monte
Carlo fatigue count replaces a derivation.

\section{Revision 11 decision and next gates}
The actual Papeeria Rev.11 already calls its diffusion an assumed closure.
Its constant-mobility FPE and interior It\^o SDE are algebraically consistent.
The first eleven displays were audited: row energy, bulk/cut counts,
fixed-geometry derivatives, affine work and density transformations can be
retained under their stated conditions. The missing links are the
thermodynamic identification and the material dynamical limits. The
abstract's wording is qualified; the original equations are preserved.

Retain Rev.11 as a conditional reference. First bind the same microscopic
coordinate, patch, measure and preparation. Then compare constrained PMF
and mean force with the rigid/relaxed energy over the occupied and escape
region. Next test conditional memory, inertial response, cross mobility,
preparation dependence and loading time scales on held-out trajectories.
Only after those gates establish compatible kinetics and boundaries can
one assign a physical clock and test an overdamped first-opening law.
An absorbing event's increasing CDF does not prove prior microstructural
damage accumulation or its equivalence to a detected crack. Later defect,
spatial and observation studies remain separate.

\appendix
\section{Source and reproducibility record}
Repository: \url{https://github.com/donghyeok8649daniel/al-fatigue-probability-theory}.
Base: \texttt{c8ccff7553cbb2e66b0c37a76ee8a24808478824}.
Files under \path{docs/research/20261020_first_principles_i_iii/}
contain the detailed six requested research notes, Korean executive summary,
scope and provenance. Isolated code is under
\texttt{research/first\_principles\_20261020/}; the immutable saved result is
\texttt{results/run\_001/verification.json}. Python 3.12.14, NumPy 2.5.3 and
SciPy 1.18.1 were used.

The downloaded Rev.11 source SHA-256 (two contiguous lines) is
\begin{quote}\ttfamily\small
ce46fad9a667ce0c6ef2e000ff9499148\\
ff7d20d1a56c427516311a9be2199be
\end{quote}
The exact source and bibliography are preserved separately. This companion
does not silently replace production potentials, calibration or historical
outputs. The full 55-display audit and actual-material fatigue validation
are outside this first delivery.

\begin{thebibliography}{9}
\bibitem{DLMF} NIST Digital Library of Mathematical Functions,
Eq.~10.32.10, \url{https://dlmf.nist.gov/10.32.E10}.
\bibitem{Mori1965} H.~Mori, Transport, collective motion, and Brownian motion,
Prog. Theor. Phys. \textbf{33}, 423--455 (1965),
\href{https://doi.org/10.1143/PTP.33.423}{doi:10.1143/PTP.33.423}.
\bibitem{Zwanzig1973} R.~Zwanzig, Nonlinear generalized Langevin equations,
J. Stat. Phys. \textbf{9}, 215--220 (1973),
\href{https://doi.org/10.1007/BF01008729}{doi:10.1007/BF01008729}.
\bibitem{Legoll2010} F.~Legoll and T.~Leli\`evre, Effective dynamics using
conditional expectations, Nonlinearity \textbf{23}, 2131--2163 (2010),
\href{https://doi.org/10.1088/0951-7715/23/9/006}{doi:10.1088/0951-7715/23/9/006}.
\bibitem{Hottovy2015} S.~Hottovy, A.~McDaniel, G.~Volpe and J.~Wehr,
The Smoluchowski--Kramers limit of stochastic differential equations with
arbitrary state-dependent friction, Commun. Math. Phys. \textbf{336},
1259--1283 (2015),
\href{https://doi.org/10.1007/s00220-014-2233-4}{doi:10.1007/s00220-014-2233-4}.
\bibitem{Ayaz2022} C.~Ayaz, B.~A.~Dalton and R.~R.~Netz, Generalized
Langevin equation with a nonlinear potential of mean force and nonlinear
memory friction from a hybrid projection scheme,
Phys. Rev. E \textbf{105}, 054138 (2022),
\href{https://doi.org/10.1103/PhysRevE.105.054138}{doi:10.1103/PhysRevE.105.054138}.
\end{thebibliography}
\end{document}
